% This is sigproc-sp.tex -FILE FOR V2.6SP OF ACM_PROC_ARTICLE-SP.CLS
% OCTOBER 2002
%
% It is an example file showing how to use the 'acm_proc_article-sp.cls' V2.6SP
% LaTeX2e document class file for Conference Proceedings submissions.
% 
%----------------------------------------------------------------------------------------------------------------
% This .tex file (and associated .cls V2.6SP) *DOES NOT* produce:
%       1) The Permission Statement
%       2) The Conference (location) Info information
%       3) The Copyright Line with ACM data
%       4) Page numbering
%
%  However, both the CopyrightYear (default to 2002) and the ACM Copyright Data
% (default to X-XXXXX-XX-X/XX/XX) can still be over-ridden by whatever the author
% inserts into the source .tex file.
% e.g.
% \CopyrightYear{2003} will cause 2003 to appear in the copyright line.
% \crdata{0-12345-67-8/90/12} will cause 0-12345-67-8/90/12 to appear in the copyright line.
%
%
%---------------------------------------------------------------------------------------------------------------
% It is an example which *does* use the .bib file (from which the .bbl file
% is produced).
% REMEMBER HOWEVER: After having produced the .bbl file,
% and prior to final submission,
% you need to 'insert'  your .bbl file into your source .tex file so as to provide
% ONE 'self-contained' source file.
%
% Questions regarding SIGS should be sent to
% Adrienne Griscti ---> griscti@acm.org
%
% Questions/suggestions regarding the guidelines, .tex and .cls files, etc. to
% Gerald Murray ---> murray@acm.org 
%
% For tracking purposes - this is V2.6SP - OCTOBER 2002


\documentclass[12pt]{article}

\setlength{\oddsidemargin}{0in}
\setlength{\evensidemargin}{0in}
\setlength{\topmargin}{0in}
\setlength{\headheight}{0in}
\setlength{\headsep}{0in}
\setlength{\textwidth}{6in}
\setlength{\textheight}{9in}
\setlength{\parindent}{0in} 

\usepackage{graphicx} %For jpg figure inclusion
\usepackage{times} %For typeface
\usepackage{epsfig}
\usepackage{xcolor} %For Comments
%\usepackage[all]{xy}
\usepackage{float}
%\usepackage{subfigure} 
\usepackage{hyperref}
\usepackage{url}
\usepackage{parskip}
\usepackage{listings}

%% Elena's favorite green (thanks, Fernando!)
\definecolor{ForestGreen}{RGB}{34,139,34}
\definecolor{BlueViolet}{RGB}{138,43,226}
\definecolor{Coquelicot}{RGB}{255, 56, 0}
\definecolor{Teal}{RGB}{2,132,130}
%Uncomment this if you want to show work-in-progress comments
\newcommand{\comment}[1]{{\bf \tt  {#1}}}
% Uncomment this if you don't want to show comments
%\newcommand{\comment}[1]{}
% for color names see https://www.overleaf.com/learn/latex/Using_colors_in_LaTeX
\newcommand{\emcomment}[1]{\textcolor{ForestGreen}{\comment{Elena: {#1}}}}
\newcommand{\todo}[1]{\textcolor{blue}{\comment{To Do: {#1}}}}
%%%%%%%%%%%%%%%%%%%%%%%%%%%%%%%%%%%%%%%%%%

\begin{document}
\pagestyle{plain}
%
% --- Author Metadata here ---
%\conferenceinfo{WOODSTOCK}{'97 El Paso, Texas USA}
%\setpagenumber{50}
%\CopyrightYear{2002} % Allows default copyright year (2002) to be
%over-ridden - IF NEED BE. 
%\crdata{0-12345-67-8/90/01}  % Allows default copyright data
%(X-XXXXX-XX-X/XX/XX) to be over-ridden. 
% --- End of Author Metadata ---

\title{15 years of Sorting Competitions.}
%\subtitle{[Extended Abstract \comment{DO WE NEED THIS?}]
%\titlenote{}}
%
% You need the command \numberofauthors to handle the "boxing"
% and alignment of the authors under the title, and to add
% a section for authors number 4 through n.
%
% Up to the first three authors are aligned under the title;
% use the \alignauthor commands below to handle those names
% and affiliations. Add names, affiliations, addresses for
% additional authors as the argument to \additionalauthors;
% these will be set for you without further effort on your
% part as the last section in the body of your article BEFORE
% References or any Appendices.

\author{
Elena Machkasova\\
Computer Science Discipline \\
University of Minnesota Morris\\
Morris, MN 56267\\
elenam@morris.umn.edu
}
\date{}
\maketitle
\thispagestyle{empty}

\section*{\centering Abstract}


\newpage
\setcounter{page}{1}
\section{Introduction}\label{sec:intro}
In the process of earning a bachelor's degree in computer science students typically learn a variety of sorting algorithms, from bubble sort to radix sort and heap sort. They implement several of these algorithms and learn their $O$ running times, advantages, and disadvantages. However, very few courses require students to implement sorting of any data types other than integers, and perhaps strings. Thus, students often graduate without actual experience of adapting algorithms to more complex data types and sorting criteria. They also don't get a clear sense for what $O$ running times really mean (for example, that $O(n \log n)$ may be more efficient than in practice than $O(n)$, depending on the constants.) 
Getting a hands-on experience with developing a fast sorting algorithm would allow students to understand the tradeoffs and the contexts of applying different algorithms, as well as the art of implementing an algorithm in a programming language, such as choosing more efficient data types and data structures.    

Recognizing the need for such an assignment, I added the \textit{Sorting Competition} in the 2010 offering of UMN Morris CSci 3501 ``Algorithms and Computability'' course.  This is a 5-credit course, covering the first two chapters (and a few selected topics) of~\cite{CLRS}, followed by the theory of computation material. The sorting competition starts after Chapter II of~\cite{CLRS} is covered and includes multiple parts:
\begin{enumerate}
\item At least two rounds of preliminary competitions
\item The final round of the competition
\item Correctness analysis
\item Presentations
\item Final reports
\end{enumerate}
The entire assignment takes about a month and usually takes place in labs while more theoretical computability material is covered in lectures. 

The competition is set up as follows: students are given a sorting task (for example, sorting positive integers by the sum of their unique prime factors). Their goal is to design and implement the fastest algorithm in Java. There are some specific restrictions which I detail in Section~\ref{subsec:java-setup}.
\emcomment{And possibly a section on the general setup, such as the timing?}
The timing uses Java {\tt time} function and is applied only to the sorting itself, not reading in the data or printing out the output. Students are given a \textit{reference implementation}: an implementation of the sorting that provides a correct order, not optimized for speed. The output of the reference implementation is compared to students' sorted data to check if the sorting is correct. This allows students to experiment with different inputs and check their correctness.  

Over the years the competition became a tradition that extends to more than just students in CSci 3501: the challenge is open to anyone interested in participating, including faculty, alumni, and students not in the class. It often features one or two ``external'' submissions, i.e. submissions from those not in the class. Some alumni participate fairly regularly. The presentations are also open to everyone. This adds a social aspects to the competition, further promoting interest in algorithmic problem solving. I am really fortunate that my colleague Dr. Peter Dolan who teaches the CSci 3501 class in some years is just as excited about the competition as I am, and we often collaborate on putting together the challenge, regardless of who of us happens to teach the class in a given year.  

In the remainder of the paper I will detail the rules and the setup in Section~\ref{sec:rules}, list some of the challenges and the takeaways in Section~\ref{sec:past}, and discuss the role of the sorting competition in Morris CSci community in Section~\ref{sec:community}. Conclusions and future directions are discussed in Section~\ref{sec:conclusions}.

\section{Rules and setup}\label{sec:rules}

\subsection{General setup and starting files}\label{subsec:general}
The goal of the competition is to come up with the fastest sorting algorithms. The speed is measured from the start of the sorting to its finish, not counting reading in the data and writing it out. This is done by using Java \texttt{System.currentTimeMillis()} method to record the times right before the sorting starts and right after it ends. The input is read from a file before the timing starts, passes the unsorted array to a timed sorting method, and the sorted array is written out to a file after the sorting is done. 
Unix \texttt{diff} command may then be used to compare the output of the students' sorting to the correct order of the data. The input/output code is provided, and there are only limited changes that can be made to it (see Section~\ref{subsec:allowed-features} for details). 

In the first couple of years of the competition students had to write their own code to check if the data is sorted correctly. However, I found out that it is more beneficial to provide a \textit{reference implementation}: an implementation of sorting that provides the correct result, but is not optimized for speed. In fact, I often provide an implementation that deliberately includes inefficiencies, such as unnecessarily repeated computations, to give students a ``low-hanging fruit'' for optimization. The reference implementation often calls a Java predefined sorting method, such as \texttt{Arrays.sort()} (TimSort since \emcomment{fill in the year}, with a comparator that implements the comparison criterion for the given data. Since many students aren't familiar with this method, it provides an additional pedagogical benefit of demonstrating the use of predefined sorting methods with a customized comparator. 

In the recent years I also started providing a data generator. There usually are data parameters that can come from specified ranges, and data often follows specific distributions. Statistical properties of data can be important, and many top-scoring algorithms adjust to specific parameters (for example, they may be different for shorter strings than for longer ones). Providing a data generator allows students to study the problem space on their own. 

An important feature of the competition is that the sorting program is ran pinned to a single core using Unix \texttt{taskset -c 0} command, so parallelization is not an optimization option. The reason for this restriction is to keep the competition within the scope of the material covered in the first two chapters of~\cite{CLRS} and to avoid giving unfair advantage to a few students who are familiar with this technique. 

\subsection{Allowed code modifications}\label{subsec:allowed-features}
The reference implementation also serves as a template for students and clearly indicates what code can and cannot be modified. Most of the students' own code must be in the timed sorting method. However, there are a few other options for modifications. There are also important constraints. Below is the list of allowed modifications and their constraints.
\begin{enumerate}
\item Depending on the problem, students may or may not be allowed to add some information gathering as the program is reading the data. However, if it is allowed,  there are limitations:  the total memory to store this information is constant w.r.t. the size of the data. Usually there is a maximum amount of memory that can be used for this purpose. This means that computing some general information, such as the minimum and the maximum value, would be fine, but recording some information about each element would not be.
\item The sorting method must take the array of elements as produced by the reading portion of the template. Typically the sorting method in the reference implementation sorts in place, and therefore is void. However, students are allowed to change its return type to return a sorted array, and its type doesn't have to be the same as the given array. In this case students must provide their own method for writing out the resulting array to a file. Recall that the method is called outside of the timed portion of the program. The constraint is that no processing is done as the data is being written out. For example, if the result needs to be in binary and of a fixed length, it's not permitted to store the result as a decimal and convert to binary with padding as the results are written out. 
\item  There are also restrictions on the use of global variables. For example, in recent years a typical restriction was to allow global constant, but not variables. There are also limits on how large the total memory for constants can be. 
\item Since the program runs on a single core, the use of \textit{multithreading} is not prohibited because it doesn't provide the benefit of parallelization.  Interestingly, in 2022 one of the teams used multithreading and discovered that it improved their time even on a single core (see Section~\ref{sec:past} for details). Thus the rules were updated to explicitly allow multithreading, but limit the number of threads to four and make it constant, not depending on the data.
\end{enumerate}

\subsection{Java setup}\label{subsec:java-setup}
Java Virtual Machine (JVM) provides a large number of automatic optimizations.  However, these optimizations are implemented in the Just-in-Time (JIT) compiler and require a method to be executed a large number of times before it is optimized. This is known as \textit{JVM warmup}~\cite{warmup}.

\subsection{Scoring and competition rounds}\label{subsec:rounds}

\section{Past challenges and takeaways}\label{sec:past}

\section{Competition as community building}\label{sec:community}

\section{Conclusions and the future}\label{sec:conclusions}

\section{Acknowledgments}\label{sec:ack}
Peter


\bibliographystyle{acm}
\bibliography{MICS2026Sorting}

\end{document}

